\section{Corrected coordinate grids}\label{sec:grids}

This section shows how a finite dynamic program yields a valid lower bound
for the continuous mixed-integer problem, and how its coordinate
min-marginals remove parts of the domain without losing any global minimizer.
Only Definition~\ref{def:curvature} is used: no convexity, uniqueness or
growth.

\subsection{Corrected grids and cell bounds}
A \emph{subbox} is a product $X'=\prod_iX'_i\subseteq X$ with
$X'_i=[\ell_i',u_i']$ for $i\in I_C$ and $X'_i=[\ell_i',u_i']\cap\Z$ with
integers $\ell_i',u_i'$ for $i\in I_Z$, where $\ell_i'\le u_i'$. A \emph{grid} for $X'$ is a product
$G=\prod_iG_i$ of finite sets $G_i\subseteq X'_i$ that contain $\ell_i'$ and
$u_i'$. The elements of $G_i$ are \emph{nodes}, and $\abs{G_i}$ is the number
of nodes per coordinate. Consecutive nodes $a<a'$ of $G_i$ form a
\emph{grid interval} $J=[a,a']$; if $\ell_i'=u_i'$, the single node $a$ forms
the grid interval $[a,a]$. The \emph{effective width} of a grid interval
$J=[a,a']$ of coordinate $i$ is $w(J)=0$ if $i\in I_Z$ and $a'-a=1$, and
$w(J)=a'-a$ otherwise; an integer interval of length one has no feasible point
in its interior.

\begin{definition}[Corrections, bound and min-marginals]\label{def:corr}
For $v\in G_i$ let $w_i(v)$ be the largest effective width of a grid interval
with endpoint $v$. Put
\begin{equation}\label{eq:correction}
d_i(v)=\frac{L_i}{8}w_i(v)^2,\qquad D(y)=\sum_{i=1}^nd_i(y_i),\qquad
Q(y)=F(y)-D(y)\qquad(y\in G),
\end{equation}
\begin{equation}\label{eq:minmarginal}
\beta=\min_{y\in G}Q(y),\qquad m_i(v)=\min\{Q(y):y\in G,\ y_i=v\}\quad(v\in G_i).
\end{equation}
The function $m_i$ is the \emph{min-marginal} of $Q$ in coordinate $i$.
\end{definition}

The corrections are unary, so $Q$ has the same scopes as $F$ and the same
tree decomposition applies to it. A \emph{cell} of $G$ is a product
$C=\prod_iJ_i$ of grid intervals $J_i$ of $G_i$; its \emph{cell bound} is
\[
\beta(C)=\min_{v\in\mathrm{vert}(C)}F(v)-\sum_i\frac{L_i}{8}w(J_i)^2,
\qquad\mathrm{vert}(C)=\prod_i\{\text{endpoints of }J_i\}.
\]
Subtracting $\frac{L_i}2(x_i-\mathrm{mid}\,J_i)^2$ in every coordinate with
$w(J_i)>0$ makes $F$ concave along each of these coordinates, so its minimum
over $C\cap X$ is attained at a vertex, and $\beta(C)$ is a lower bound for
$F$ on $C\cap X$. For a single cell
this is the vertex bound of Bajaj and Hasan \cite{BajajHasan2020}, who build
an edge-concave underestimator from an upper bound on the diagonal Hessian
entries and evaluate it at the vertices; we use one bound $L_i$ per
coordinate. The term $\frac{L_i}8w(J_i)^2$ is the maximum
separation of the $\alpha$BB underestimator with $\alpha_i=L_i/2$
\cite{MaranasFloudas1994,AdjimanEtAl1998}, used with the opposite sign. The
following observation shows that the unary correction~\eqref{eq:correction}
computes the best of these bounds over all cells of a product grid at once.

\begin{proposition}[Cellwise form of the corrected grid]\label{prop:cellwise}
Let $G$ be a grid for a subbox $X'$.
\begin{enumerate}[label=(\alph*)]
\item For every cell $C$ and every $x\in C\cap X$, $F(x)\ge\beta(C)$. Every
$x\in X'$ lies in some cell.
\item $\beta=\min_C\beta(C)$, the minimum over all cells. In particular
$\beta\le\min_{X'}F$.
\item For a grid interval $J=[a,a']$ of coordinate $i$, put
$\beta_i(J)=\min\{\beta(C):C_i=J\}$. Then
\[
\beta_i(J)=\min_{v\in\{a,a'\}}\bigl(m_i(v)+d_i(v)\bigr)-\frac{L_i}8w(J)^2
\ \ge\ \min\{m_i(a),m_i(a')\},
\]
and $F(x)\ge\beta_i(J)$ for every $x\in X'$ with $x_i\in J$. Moreover,
$F(x)\ge m_i(v)+d_i(v)$ for every $x\in X'$ with $x_i=v\in G_i$.
\end{enumerate}
\end{proposition}

\begin{proof}
(a) Let $C=\prod_iJ_i$ with $J_i=[a_i,a_i']$, let $W=\{i:w(J_i)>0\}$,
$c_i=(a_i+a_i')/2$, and
$\varphi(y)=F(y)-\sum_{i\in W}\frac{L_i}{2}(y_i-c_i)^2\le F(y)$. If $i\notin W$,
then $a_i'=a_i$, or $i\in I_Z$ and $a_i'=a_i+1$; in both cases
$x_i\in\{a_i,a_i'\}$. If $i\in W$, then for fixed values of the other
coordinates,
$t\mapsto\varphi(\dots,t,\dots)=\bigl(F(\dots,t,\dots)-\frac{L_i}2t^2\bigr)
+L_ic_it+\mathrm{const}$ is concave on $J_i$ by
Definition~\ref{def:curvature}, so its minimum over $J_i$ is attained at an
endpoint. Replacing the coordinates in $W$ one at a time by such endpoints
keeps the point in $C\cap X$ and does not increase $\varphi$, and ends at a
vertex $v$ with $\varphi(v)\le\varphi(x)$. Since $(v_i-c_i)^2=w(J_i)^2/4$
for $i\in W$,
$F(x)\ge\varphi(x)\ge\varphi(v)=F(v)-\sum_i\frac{L_i}8w(J_i)^2\ge\beta(C)$.
Every $x_i\in X'_i$ lies in a grid interval because $G_i$ contains both
endpoints of $X'_i$.

(b) For $v\in G_i$ let $\mathcal J_i(v)$ be the nonempty set of grid intervals
with endpoint $v$. The pairs $(C,y)$ with $y\in\mathrm{vert}(C)$ are exactly
the grid points $y$ together with an independent choice of
$J_i\in\mathcal J_i(y_i)$ for every $i$. Hence
\[
\min_C\beta(C)=\min_{y\in G}\Bigl[F(y)-\sum_i\max_{J\in\mathcal J_i(y_i)}
\frac{L_i}8w(J)^2\Bigr]=\min_{y\in G}\bigl[F(y)-D(y)\bigr]=\beta ,
\]
and $\beta\le\min_{X'}F$ follows from (a).

(c) The same computation with the $i$th interval fixed to $J$ gives the
formula, because $Q(y)=F(y)-\sum_{k\ne i}d_k(y_k)-d_i(v)$ when $y_i=v$. Since
$J\in\mathcal J_i(v)$ for $v\in\{a,a'\}$, $d_i(v)\ge L_iw(J)^2/8$, which
gives the inequality. The bound for $x_i\in J$ follows from (a) with $C_i=J$.
If $x_i=v\in G_i$, choose $J_i\in\mathcal J_i(v)$ and grid intervals
$J_k\ni x_k$ for $k\ne i$, and run the argument of (a) with $W$ replaced by
$W\setminus\{i\}$. Coordinate $i$ stays at $v$, and the vertex $y$ reached
satisfies
$F(x)\ge F(y)-\sum_{k\ne i}\frac{L_k}8w(J_k)^2\ge F(y)-\sum_{k\ne i}d_k(y_k)
=Q(y)+d_i(v)\ge m_i(v)+d_i(v)$.
\end{proof}

Enumerating all cells and their up to $2^n$ vertices is the box-by-box
computation of branch-and-bound. Because $D$ is a sum of unary terms, the same
number $\beta$ is the minimum of a factorized function on the product grid,
and Section~\ref{sec:grids-dp} computes it, with all min-marginals, by
dynamic programming on the given tree decomposition.

\paragraph{Families of intervals.}
The proof of Proposition~\ref{prop:cellwise} uses the grid intervals only
through two properties: each is an interval with endpoints in $G_i$, and
every node is an endpoint of at least one of them. The proposition and its
proof therefore remain valid for finite sets $G_i\subseteq X_i$ and, in place of the grid
intervals, finite families $\mathcal J_i$ of intervals with endpoints in
$G_i$ such that every node is an endpoint of a member, with $w_i(v)$ the
largest effective width of a member with endpoint $v$, cells the products of
members, and $X'$ replaced by $X\cap\prod_i\bigcup\mathcal J_i$. The
bounds of Lemma~\ref{lem:cells} for unions of uniform cells are an instance
of this form.

\paragraph{Randomized rounding.}
The bound $\beta\le F(x)$ also follows from randomized rounding. For
$x\in X'$ choose grid intervals $J_i\ni x_i$ and round each coordinate
independently to an endpoint of $J_i$ with mean $x_i$; call the resulting
grid point $Y$. Jensen's inequality for the concave functions
$t\mapsto F(\dots,t,\dots)-\frac{L_i}2t^2$, applied one coordinate at a time,
gives $\E F(Y)\le F(x)+\sum_i\frac{L_i}2\Var(Y_i)$. Since
$\Var(Y_i)\le w(J_i)^2/4$ (and $Y_i=x_i$ if $w(J_i)=0$) and
$d_i(Y_i)\ge L_iw(J_i)^2/8$, this gives $\E Q(Y)\le F(x)$.
Section~\ref{sec:constraints} uses a correlated version of this rounding that
keeps coupling constraints satisfied (Lemma~\ref{lem:tu-round}).

\subsection{Dynamic programming and min-marginals}\label{sec:grids-dp}
Let $q_t$ be the sum of the factors assigned to bag $t$ minus the corrections
$d_i$ of the coordinates whose home bag is $t$, regarded as a table on
$G_{B_t}=\prod_{i\in B_t}G_i$; its entries are the \emph{table entries} of bag
$t$, and $Q=\sum_tq_t$ on $G$. For adjacent tree nodes $t,u$ let
$S_{tu}=B_t\cap B_u$ and define messages
\begin{equation}\label{eq:messages}
M_{t\to u}(y_{S_{tu}})=\min_{y_{B_t\setminus S_{tu}}}\Bigl[q_t(y_{B_t})
+\sum_{v\in\mathcal N(t)\setminus\{u\}}M_{v\to t}(y_{S_{vt}})\Bigr],
\end{equation}
where $\mathcal N(t)$ is the set of neighbors of $t$ in $T$ and an empty
separator carries a single scalar. Put
$A_t=q_t+\sum_{v\in\mathcal N(t)}M_{v\to t}$.

\begin{lemma}[Two-pass dynamic programming]\label{lem:dp}
For every bag $t$ and every $y_{B_t}\in G_{B_t}$,
$A_t(y_{B_t})=\min\{Q(z):z\in G,\ z_{B_t}=y_{B_t}\}$. Hence $\beta=\min A_t$, and
$m_i(v)$ is the minimum of $A_t$ over $y_i=v$ for any bag $t\ni i$. With
$K=\max_i\abs{G_i}$, all messages, $\beta$, a minimizer and all min-marginals
are obtained with $O\bigl((\abs{\mathcal A}+n+N)K^p\bigr)$ additions and
comparisons of table entries and $O(p)$ index work per table entry, hence
$O\bigl(p(\abs{\mathcal A}+n+N)K^p\bigr)$ operations, plus the evaluation of
every factor at the at most $K^{\abs{S_a}}$ grid points of its scope.
\end{lemma}

\begin{proof}
For an oriented edge $t\to u$ let $T_{t\to u}$ be the component of
$T\setminus\{tu\}$ containing $t$, and $\mathcal W_{t\to u}$ the set of
coordinates occurring in bags of $T_{t\to u}$ but not in $S_{tu}$. By
induction on $\abs{T_{t\to u}}$, $M_{t\to u}(y_{S_{tu}})$ is the minimum of
$\sum_{r\in T_{t\to u}}q_r$ over $y_{\mathcal W_{t\to u}}$: the trees
$T_{v\to t}$, $v\in\mathcal N(t)\setminus\{u\}$, partition
$T_{t\to u}\setminus\{t\}$, and by running intersection a coordinate occurring
both in $T_{v\to t}$ and outside it lies in $S_{vt}$, so the sets
$\mathcal W_{v\to t}$ are disjoint from each other and from $B_t$ and the
minimization splits. The same argument with all neighbors of $t$ gives the
formula for $A_t$. Messages are computed toward a root and back; at bag $t$
the sum of all incoming messages is formed once, and each outgoing message is
obtained from $A_t-M_{u\to t}$ by one pass, with no product over children. A
minimizer is recovered from the root outward.
\end{proof}

Lemma~\ref{lem:dp} is classical nonserial dynamic programming, bucket
elimination or junction-tree computation \cite{BerteleBrioschi1972,Dechter1999,
LauritzenSpiegelhalter1988}, and two passes give all min-marginals
\cite{WainwrightJaakkolaWillsky2005,KohliTorr2006}. The only new element is
that its input is the corrected objective $Q$. Any exact method for $\beta$, a
minimizer and the min-marginals can replace the messages;
Section~\ref{sec:balanced} uses minimum cuts for balanced quadratics without
a width bound.

\subsection{Filtering and certificates}

\begin{definition}[Filtering step]\label{def:filter}
Let $G$ be a grid for a subbox $X'$ with corrected minimum $\beta$, and let
$U\ge\beta$. For $i\in P$, \emph{retain} a grid interval $[a,a']$ of
coordinate $i$ if
\begin{equation}\label{eq:retain}
\min\{m_i(a),m_i(a')\}\le U ,
\end{equation}
and let $X_i''$ be the smallest interval containing the retained intervals
(intersected with $\Z$ for $i\in I_Z$). For $i\notin P$ let $X_i''=X_i'$. The
filtered subbox is $X''=\prod_iX_i''$.
\end{definition}

\begin{proposition}[Safe filtering]\label{prop:filter}
Every $x\in X'\setminus X''$ satisfies $F(x)>U$. Hence $X''$ contains every
$x\in X'$ with $F(x)\le U$, and every minimizer of $F$ over $X'$ if
$U\ge\min_{X'}F$. Every minimizer $y$ of $Q$ over $G$ lies in $X''$; in
particular $X''$ is a subbox.
\end{proposition}

\begin{proof}
If $x\in X'\setminus X''$, some $i\in P$ has $x_i\notin X_i''$, so every grid
interval containing $x_i$ was not retained, and
Proposition~\ref{prop:cellwise}(c) gives $F(x)\ge\min\{m_i(a),m_i(a')\}>U$
for such an interval. For a minimizer $y$ of $Q$ and $i\in P$,
$m_i(y_i)\le Q(y)=\beta\le U$, and $y_i$ is an endpoint of a grid interval of
coordinate $i$, which is therefore retained. So $y_i\in X_i''$ for $i\in P$,
and trivially for $i\notin P$. Hence each $X_i''$ is nonempty, and for
$i\in I_Z$ its endpoints are nodes, hence integers.
\end{proof}

A method that filters repeatedly produces nested subboxes. The following
certificate records only what is needed to check the result.

\begin{definition}[Path certificate]\label{def:cert}
A \emph{path certificate} consists of grids $G^{(0)},\dots,G^{(k)}$, a number
$\beta$ and a point $\hat x$. Let $X^{(j)}$ be the subbox spanned by
$G^{(j)}$: with $\ell_i^{(j)}=\min G_i^{(j)}$ and $u_i^{(j)}=\max G_i^{(j)}$,
$X^{(j)}_i$ is $[\ell_i^{(j)},u_i^{(j)}]$ or its integer points. The
certificate is \emph{valid} if $\hat x\in X$, the nodes of integer coordinates
are integers, $X^{(0)}=X$, $X^{(j+1)}\subseteq X^{(j)}$, and, with $Q^{(j)}$
and $m^{(j)}$ computed from $G^{(j)}$ and the curvature bounds,
\begin{enumerate}[label=(C\arabic*)]
\item for $j<k$, every grid interval $J=[a,a']$ of $G^{(j)}_i$ that is not
contained in $[\ell_i^{(j+1)},u_i^{(j+1)}]$ satisfies
\[
\min\bigl\{m^{(j)}_i(a),m^{(j)}_i(a')\bigr\}\ge\beta ,
\]
or has $w(J)=0$ and $m^{(j)}_i(v)\ge\beta$ for each endpoint $v$ of $J$
outside $[\ell_i^{(j+1)},u_i^{(j+1)}]$;
\item $\min_{G^{(k)}}Q^{(k)}\ge\beta$.
\end{enumerate}
\end{definition}

The second alternative in (C1) covers filters that remove single integer
nodes, such as the integer-label filter of Lemma~\ref{lem:labelfilter}.

\begin{theorem}[Path certificates]\label{thm:certificate}
If a path certificate is valid, then $\beta\le\OPT\le F(\hat x)$. Its
validity is checked by at most $k+1$ runs of the dynamic program of
Lemma~\ref{lem:dp} in exact arithmetic, given the curvature bounds; for a
quadratic the checker computes $L_i=\max\{H_{ii},0\}$ itself. No growth or
uniqueness assumption is involved.
\end{theorem}

\begin{proof}
Let $x\in X$ and let $j$ be the largest index with $x\in X^{(j)}$. If $j=k$,
Proposition~\ref{prop:cellwise}(b) and (C2) give $F(x)\ge\beta$. Otherwise
some $x_i\notin X^{(j+1)}_i$, so $x_i\notin[\ell_i^{(j+1)},u_i^{(j+1)}]$,
because $x_i$ is an integer if $i\in I_Z$. A grid interval $J$ of $G_i^{(j)}$
containing $x_i$ is therefore not contained in
$[\ell_i^{(j+1)},u_i^{(j+1)}]$. If the first alternative of (C1) holds for
$J$, Proposition~\ref{prop:cellwise}(c) gives $F(x)\ge\beta$. Otherwise
$w(J)=0$, so the feasible value $x_i$ is an endpoint of $J$ outside
$[\ell_i^{(j+1)},u_i^{(j+1)}]$, and the last claim of
Proposition~\ref{prop:cellwise}(c) gives
$F(x)\ge m^{(j)}_i(x_i)+d^{(j)}_i(x_i)\ge m^{(j)}_i(x_i)\ge\beta$.
\end{proof}

If $G^{(0)}$ spans $X$, each grid $G^{(j+1)}$ spans the subbox obtained from
$G^{(j)}$ by the filtering step with a threshold $U_j$ that is the value of a feasible point,
and $\beta$ is the corrected minimum of the last grid, then the record is a
valid path certificate. Indeed, by Proposition~\ref{prop:filter} and
induction every minimizer lies in every $X^{(j)}$, so the corrected minimum
$\beta_j$ of $G^{(j)}$ satisfies $\beta_j\le\OPT\le U_j$, as
Definition~\ref{def:filter} requires; in particular $\beta\le\OPT$. A grid
interval of $G_i^{(j)}$ not contained in $[\ell_i^{(j+1)},u_i^{(j+1)}]$ was
not retained, so its endpoint bound exceeds $U_j\ge\beta$. Thus a certificate
consists of the stage grids (or the boxes with the centers, meshes and
grading that generate them), one bound and one feasible point. Messages,
min-marginal tables and intermediate incumbents need not be stored, and stages
other than the last that remove nothing may be dropped. The stages that remove
points are needed: the last grid bounds $F$ only on $X^{(k)}$. A checker runs
the dynamic program once per recorded stage, the same table work as the
filtering steps that produced the certificate;
Section~\ref{sec:computation} reports measured replay times.

The certificate is a branch-and-bound proof whose tree is a path. Removing an
interval whose bound exceeds an incumbent is optimality-based bound
tightening, and recording these decisions follows checkable branch-and-bound
certificates (Section~\ref{sec:related}). What is specific here is that all
interval bounds of a stage come from one pair of message passes.
